


\setlength{\columnsep}{15pt}%
\setlength{\intextsep}{0pt plus 0pt minus 0pt}

\setlength{\parindent}{\myindent}
\setlength{\parskip}{\myparskip}

 \vspace*{-20pt}

% Artificial intelligence now informs consequential decisions in domains like clinical diagnosis, content moderation, and scientific peer review.
% In these settings, we expect human and AI to make better decisions together, e.g., models may err as a result of problem or distribution misspecification, experts may have access to hard-to-encode information that models do not, and accountability ultimately rests with human decision-makers. 
% The promise of AI is therefore \emph{complementarity}, i.e., collaborative decisions that outperform what a single agent (human or machine) could achieve alone.
% However, a decade of empirical work delivers largely negative evidence: human-AI combinations frequently perform worse than the better of the two alone, and reliable interventions for improving decision quality remain elusive.
% As increasingly capable AI agents are deployed alongside human experts and other agents, whether and how they add value to the decisions of existing agents is among the most consequential open problems in the field.

Artificial intelligence now informs consequential decisions in domains like clinical diagnosis, content moderation, and scientific peer review.
In these settings, we expect human and AI to make better decisions together, e.g., models may err from problem misspecification, experts may hold hard-to-encode information that models lack, and accountability ultimately rests with human decision-makers. 
The promise of AI is therefore \emph{complementarity}, i.e., collaborative decisions that outperform either agent alone.
However, a decade of empirical work delivers largely negative evidence: human-AI combinations frequently perform worse than the better of the two alone, and reliable interventions remain elusive.
As increasingly capable AI agents are deployed alongside human experts, whether and how they add value to existing agents' decisions is among the field's most consequential open problems.

% \textbf{I use decision theory to formalize and design algorithms to achieve the goal of complementarity in multi-agent collaboration. } 
% % My research develops formal estimands and diagnostic frameworks~\cite{guo2024decision, guo2026value, guo2026explanations} and applicable algorithms~\cite{guo2026complllm}, with human behavioral experiments as external validation~\cite{guo2026value, guo2026complllm}.
% Instead of relying on intuitions, I first quantify the theoretical possible complementary value with formal estimands and diagnostic frameworks~\cite{guo2024decision, guo2026value, guo2026explanations}, and then develop interventions grounded in the theory~\cite{guo2026complllm}, with behavioral experiments as external validation~\cite{guo2026value, guo2026complllm}.

\textbf{I use decision theory to formalize and design algorithms to achieve the goal of complementarity in multi-agent collaboration. } 
Rather than relying on intuition, I first quantify the theoretically possible complementary value with formal estimands and diagnostic frameworks~\cite{guo2024decision, guo2026value, guo2026explanations}, then develop interventions grounded in theory~\cite{guo2026complllm}, with behavioral experiments as external validation~\cite{guo2026value}.

\textbf{A defining characteristic of my research is its systematic interdisciplinary approach.}
Human-AI collaboration has been studied in parallel by three communities, and my work combines their strengths: rigorous formalization from theory and economics; modern post-training methods from machine learning; and controlled behavioral studies from human-computer interaction.
This combination lets me answer questions no single community's toolkit can---not whether a model is accurate, but whether it has the potential to improve decisions, and whether the potential translates to actual performance.

% Human-AI collaboration has been studied in parallel by three communities, .
% My research combines the strengths of three communities who have studied the problem of human-AI collaboration in parallel: theory and economics, machine learning community  and human-computer interaction. 
% From theory, I bring rigorous formalization that turns loose notions of ``reliance'' and ``complementarity'' into measurable estimands with rational-agent upper bounds. 
% From machine learning, I bring realistic tasks and modern post-training methods that optimize LLMs for complementary value directly. 
% From human-computer interaction, I bring controlled behavioral studies with human participants and LLM decision-makers that test whether theoretically available value is actually realized. 
% This combination lets me answer questions that no single community's toolkit can--not just whether a model is accurate, but whether its information has the potential and also really improves the decisions of agents.

My work has appeared at FAccT 2024, ICLR 2026, and ICML 2026, with ComplLLM accepted at NeurIPS 2026, and has been recognized with invited talks at the 2025 NSF AI-SDM workshop, the 2026 ESIF Economics and AI+ML Meeting, and the 2026 INFORMS Annual Meeting.
Looking ahead, this interdisciplinary skill set positions me for a future agenda to establish complementarity as a principled target for AI development, where AIs are trained, benchmarked, and deployed not to just match but to measurably improve the human experts.
